\documentclass[11pt,a4paper]{article}
\usepackage[margin=1.8cm]{geometry}
\usepackage{amsmath,amssymb,tabularx,array,needspace,fontspec,enumitem,cancel}
\newcommand{\tb}{\cancel{||||}\,}
\setmainfont{FreeSerif}
\newfontfamily\urdufont[Script=Arabic]{FreeSerif}
\TeXXeTstate=1
\newcommand{\urdu}[1]{{\urdufont\beginR #1\endR}}
\setlength{\parindent}{0pt}
\setlength{\parskip}{0pt}
\renewcommand{\arraystretch}{1.15}
\pagestyle{plain}\hyphenpenalty=10000 \exhyphenpenalty=10000

% \mcq{number}{question}{A}{B}{C}{D}
\newcommand{\mcq}[6]{%
  \Needspace{4\baselineskip}%
  \par\medskip\noindent\textbf{#1.}\ #2\par\nopagebreak
  \noindent\begin{tabularx}{\linewidth}{@{\hspace{1.5em}}*{4}{>{\raggedright\arraybackslash}X}@{}}
    (A)~#3 & (B)~#4 & (C)~#5 & (D)~#6
  \end{tabularx}\par}

\begin{document}

\begin{center}
  {\Huge\bfseries The Hope Science Academy}\\[3pt]
  {\small Excellence is our standard not goal}\\[2pt]
  {\small 247-E-1, (Opposite LDA School) Johar Town, Lahore. PH: 0300-8194789, 0322-4157001}
\end{center}

\medskip
\noindent
\begin{tabularx}{\linewidth}{|X|X|X|X|}
\hline
\textbf{Student Name:} & \textbf{Roll Num:} & \textbf{Class:} 11th & \textbf{Paper Code:} \\ \hline
\textbf{Subject Name:} Principles of Accounting & \textbf{Time Allowed:} & \textbf{Total Marks:} 405 & \textbf{Exam Date:} 02-Nov-2026 \\ \hline
\multicolumn{4}{|l|}{\textbf{Exam Syllabus:} FULL BOOK} \\ \hline
\end{tabularx}

\bigskip
\noindent\textbf{Q1. Choose the correct answer.}\hfill\textbf{1$\times$405=405}

\mcq{1}{The prime function of accounting is to:}{Record economic}{Classifying and recording business transactions}{Classifying and summarizing business transactions}{Provide the informational basis for action}
\mcq{2}{Any legal activity which is done for the purpose of earning profit, is called:}{Trust}{Society}{None of “A”, “B”}{Business}
\mcq{3}{All those things which are purchased for resale purpose are called:}{Goods}{Assets}{Liabilities}{Sale}
\mcq{4}{Book-keeping provides the:}{Primary function}{Final information}{Secondary function}{All of the above}
\mcq{5}{Book-keeping is mainly concerned with:}{Recording of business transactions}{Classifying of business transactions}{Interpreting of business transactions in money terms}{None of the above}
\mcq{6}{How many branches of accounting are:}{Two}{Three}{Four}{One}
\mcq{7}{Accounting is the language of:}{Business}{School}{America}{Proprietor}
\mcq{8}{The rate of allowance given by a creditor to a debtor, if he pays his debts before the due date of payment is normally:}{2\%}{5\%}{10\%}{15\%}
\mcq{9}{All those goods which are lying unsold in a business are termed as:}{Revenue}{Expense}{Stock}{Sale}
\mcq{10}{The term expenses and expenditure are:}{Same}{Different}{Opposite}{Singular or plural}
\mcq{11}{A person who owes money to the business is called:}{Debtor}{Creditor}{Shareholder}{Owner}
\mcq{12}{Stages for proceeding a transaction are:}{Recording}{Classifying}{Summarizing}{All of these}
\mcq{13}{If damaged or below standard goods returned to the seller are called:}{..........Returns outwards}{Purchases returns}{Returns to supplier}{All of these}
\mcq{14}{The system of providing quantified information about an organization to people who need such information is called:}{Economics}{Accounting}{Book Keeping}{None of above}
\mcq{15}{The amount of cash or goods invested by the proprietor in a business is called:}{Revenue}{Asset}{Capital}{Expense}
\mcq{16}{A person from whom credit purchase are made is called:}{Debtor}{Creditor}{Banker}{Owner}
\mcq{17}{The price of goods sold or services provided by a business to its customers is called:}{Expense}{Lose}{Liability}{Revenue}
\mcq{18}{A form of remuneration for services rendered by one person to another:}{Discount allowed}{Discount received}{Commission}{Expenditure}
\mcq{19}{It takes place when assets or services are acquired:}{Revenue}{Expenses}{Asset}{Expenditure}
\mcq{20}{The discount allowed by wholesaler to the retailer at the list price of goods is called:}{Trade discount}{Cash discount}{Discount allowed}{Discount received}
\mcq{21}{A person to whom the goods or services are sold on credit basis is known as:}{Creditor}{Proprietor}{Debtor}{Investor}
\mcq{22}{Net income is equal to:}{Assets-Liabilities}{Revenues+Expense}{Revenue-Income}{Revenues- Expenses}
\mcq{23}{The system of recording transactions having two fundamental aspects is known as:}{Single entry system}{Modern System}{Double entry system}{None of these}
\mcq{24}{According to this concept, it is assumed that business will exit for an indefinite time period:}{Realization Concept}{Going concern concept}{Business entity concept}{None of these}
\mcq{25}{According to this concept, business and owner both have separate identity:}{Going concern concept}{Business entity concept}{Cost concept}{Realization concept}
\mcq{26}{According to this concept, a record is made only of those information that can be expressed in monetary terms:}{Non-monetary concept}{Money measurement concept}{Going concern concept}{Accounting period concept}
\mcq{27}{According to this concept, an asset is recorded at a price at which it is acquired:}{Money measurement concept}{Dual aspect concept}{Cost concept}{Realization concept}
\mcq{28}{Modern accounting is based on:}{Matching concept}{Dual aspect concept}{Cost concept}{Going concern concept}
\mcq{29}{According to this concept, expenses are matched with revenues to study the business result:}{Matching concept}{Dual aspect concept}{Cost concept}{Business entity concept}
\mcq{30}{According to this concept, revenue should be recognized at the time when goods are sold or services are rendered:}{Money measurement concept}{Dual aspect concept}{Cost concept}{Realization concept}
\mcq{31}{According to this convention, anticipate no profits but provide for all possible losses:}{Full disclosure}{Conservation}{Consistency}{Materiality}
\mcq{32}{According to this convention, the users of financial statement are informed of any facts necessary for the proper interpretation of statements:}{Conservatism}{Materiality}{Full disclosure}{Consistency}
\mcq{33}{According to this convention, accounting practice should remain unchanged from one period to another:}{Conservation}{Materiality}{Full disclosure}{Consistency}
\mcq{34}{A type of business in which a merchant purchases goods and sells it in the same position is called:}{Trading concern}{Service concern}{Manufacturing concern}{None of these}
\mcq{35}{When there is no question of meeting the value of a transaction, it will be known as:}{Cash Transaction}{Credit Transaction}{Paper Transaction}{None of these}
\mcq{36}{If there is no change in total value of assets and liabilities of a business concern it will be considered as:}{Qualitative change}{Quantitative change}{Monetary change}{Economic change}
\mcq{37}{According to accounting equation capital is equal to:}{Assets+Liabilites}{Expenses- Income}{Liabilities-Assets}{Assets-Liabilities}
\mcq{38}{Expenses paid by a business decrease:}{Cash}{Capital}{Cash \& Capital}{None of these}
\mcq{39}{Every transaction affects:}{One item}{Two items}{Three items}{So many items}
\mcq{40}{A person who owns money to a business is known as a:}{Creditor}{Debtor}{Investor}{Solvent}
\mcq{41}{If furniture purchased on credit basis, it will be:}{Increase in assets and increase in capital}{Increase in assets and increases in expenses}{Increase in assets and increase in equity}{Increase in assets and increase in liabilities}
\mcq{42}{Resources owned by the business:}{Assets}{Liabilities}{Capital}{Profit}
\mcq{43}{Furniture purchased or domestic use:}{Increase in assets and increase in expenses}{Increase in expenses and increase in liabilities}{Decrease in capital and decrease in assets}{Increase in assets and decrease in assets}
\mcq{44}{When cash is paid to debtors, it will decrease:}{Cash}{Capital}{Debtor}{None of these}
\mcq{45}{Claims against assets owned by business are called:}{Assets}{Liabilities}{Capital}{Equities}
\mcq{46}{A transaction taking place with an outside person or organization is called:}{Internal Transaction}{External Transaction}{Paper Transaction}{None of these}
\mcq{47}{Cash received from debtors, it will be:}{Increase in liabilities and decrease in assets}{Increase in assets and decrease in revenue}{Increase in assets and decrease in assets}{None of these}
\mcq{48}{The accounting equation represents:}{Resources in the business are equal to the sources of business}{Resources are allocated at cost price}{Resources in the business are not equal to the sources of business}{Owner’s give money for business}
\mcq{49}{The excess of assets over liabilities is called:}{Profit}{Income}{Capital}{Equities}
\mcq{50}{Sold goods costing Rs1500 for Rs 1750 will increase the owner’s equity by Rs:}{250}{1500}{1750}{3250}
\mcq{51}{Purchased goods on credit and for cash will affect:}{Cash and goods}{Cash and creditors}{Cash, creditors and owner’s equity}{Cash, goods and creditors}
\mcq{52}{Received cash from Azhar will affect:}{Cash and onwer’s equity}{Cash and creditors}{Cash and debtors}{Cash and goods}
\mcq{53}{Cash paid to Nadeem will affect:}{Cash and owner’s equity}{Cash and creditors}{Cash and debtors}{Cash and goods}
\mcq{54}{All events, which are related to business, are measured in:}{Kilograms}{Money}{Quality}{Quantity}
\mcq{55}{The rights possessed by the owner against the assets of the business are called:}{Assets}{Liabilities}{Owner’s Equity}{Owner}
\mcq{56}{All debts which are payable by a business to its investors and others are known as:}{Capital}{Equities}{Expenses}{Liabilities}
\mcq{57}{The best system of accounting in the modern world is:}{Single entry system}{Double entry system}{Multiple entry system}{No entry system}
\mcq{58}{Discount allowed is:}{An asset}{A revenue}{A liability}{An expense}
\mcq{59}{Accounts which keep records of properties or things are called:}{Personal accounts}{Nominal accounts}{Real accounts}{None of these}
\mcq{60}{Accounts which keep records of expenses, gains and losses are called:}{Personal accounts}{Nominal accounts}{Real accounts}{None of these}
\mcq{61}{Purchase of furniture should be debited to:}{Purchase account}{Creditor account}{Furniture account}{Cash account}
\mcq{62}{Rent, salaries and wages paid are:}{Nominal accounts}{Real accounts}{Personal accounts}{None of these}
\mcq{63}{The claim of the supplier in the business is called:}{Revenue}{Liability}{Expense}{Asset}
\mcq{64}{Ahmad, a account is an example of:}{Personal account}{Real account}{Nominal account}{None of these}
\mcq{65}{How many categories of account are:}{Two}{Three}{Four}{Five}
\mcq{66}{Goods purchased for cash from Babar should be debited to:}{Babar account}{Goods account}{Purchase account}{Cash account}
\mcq{67}{Stationery account is an example of:}{Expense}{Loss}{Consumable asset}{None of these}
\mcq{68}{Discount received from Ali should be debited to:}{Discount received account}{Cash account}{Ali account}{None of these}
\mcq{69}{Cash invested in business by owner will increase:}{Cash}{Capital}{Liability}{Both a,b}
\mcq{70}{Discount allowed to Ahmad will be credited to:}{Ahmad’s account}{Discount allowed account}{Cash account}{None of these}
\mcq{71}{Cash and personal accounts are maintained in:}{Single entry system}{Double entry system}{Multiple entry system}{None of these}
\mcq{72}{Real and Nominal accounts are maintained under:}{Single entry system}{Double entry system}{Triple entry system}{Multiple entry system}
\mcq{73}{Rent paid to landlord for a portion of building will be debited to:}{Landlord}{Building account}{Rent account}{Cash account}
\mcq{74}{Purchased goods on credit increases the:}{Creditors}{Debtors}{Customer}{None of these}
\mcq{75}{Purchased furniture on account will be credited to:}{Purchased account}{Furniture account}{Cash account}{Creditor account}
\mcq{76}{All expense and losses should be debited and all gains should be credited is applied to:}{Personal accounts}{Real accounts}{Nominal accounts}{None of these}
\mcq{77}{Assets coming in should be debited and going out should be credited is applied to:}{Personal accounts}{Real accounts}{Nominal accounts}{A and B}
\mcq{78}{A person who receive the benefit should be debited and who give the benefit should be credited is applied to:}{Personal accounts}{Real accounts}{Nominal accounts}{B and A}
\mcq{79}{Debit the increase in:}{Assets, Expenses}{Assets, Revenue}{Assets, Liabilities}{Assets, Capital}
\mcq{80}{Credit the increase in:}{Asset, Expenses}{Assets, Capital and Liabilities}{Revenue, Capital and Liabilities}{Expenses, Capital and Liabilities}
\mcq{81}{Debit the decrease in:}{Asset, Expenses}{Revenue, Capital and Liabilities}{Assets, Capital and Liabilities}{Expenses, Capital and Liabilities}
\mcq{82}{Credit the decrease in:}{Revenue and Liabilities}{Expenses and Liabilities}{Assets and Liabilities}{Asset and Expenses}
\mcq{83}{Sold merchandises on credit to Ali Rs7000 involves:}{Nominal and Real accounts}{Real and Personal accounts}{Nominal and personal accounts}{Real and Real accounts}
\mcq{84}{Invested Rs2,00,000 Cash, Building Rs75000 and furniture Rs 25,000 involves:}{Real and personal accounts}{Real and Nominal accounts}{Personal and Nominal accounts}{Real accounts}
\mcq{85}{Purchased Goods for Cash Rs 50000 subject to 10\% trade discount involves:}{Real and Personal accounts}{Real and Nominal accounts}{Personal and Nominal accounts}{Nominal accounts}
\mcq{86}{Withdrew cash for personal use Rs300 involves:}{Real and Proprietorship account}{Proprietorship and Personal accounts}{Real and Nominal accounts}{Nominal and Personal accounts}
\mcq{87}{Discount received is a:}{Asset}{Revenue}{Liability}{Expense}
\mcq{88}{A summary of transactions relating to a person or thing:}{Accounting}{Book-keeping}{Account}{Accountancy}
\mcq{89}{Journal is called the book of:}{Original entry}{First entry}{Preliminary entry}{All of above}
\mcq{90}{The entry in which more than one account is debited or credited is called:}{Simple entry}{Compound entry}{Single entry}{Prime entry}
\mcq{91}{The art of recording transaction in a journal is called:}{Posting}{Journaling}{Ledger entry}{Compound entry}
\mcq{92}{Recording a transaction in the appropriate place of the concerned book of account is called:}{A statement}{Business document}{An entry}{None of these}
\mcq{93}{A short explanation of each transaction is written under each entry is called:}{Posting}{Recording}{Summarizing}{Narration}
\mcq{94}{The money with which proprietor starts his business is called:}{Cash}{Investment}{Capital}{Aid}
\mcq{95}{Goods purchased on credit from Tariq should be debited to:}{Goods account}{Purchase account}{Tariq account}{Cash account}
\mcq{96}{Purchased a motorcycle by the proprietor for his personal use should be debited to:}{Purchases account}{Motorcycle account}{Proprietor’s Personal account}{Drawings account}
\mcq{97}{Goods sold on account to Babar should be credited to:}{Babar account}{Goods account}{Sales account}{Purchase account}
\mcq{98}{Goods returned to Ahmad should be debited to:}{Sales return account}{Ahmad account}{Purchase return account}{Purchase account}
\mcq{99}{Goods returned by Zain should be debited to:}{Purchase return account}{Zain account}{Sales return account}{Purchase account}
\mcq{100}{When goods are given away as charity or free samples should be credited to:}{Charity account}{Free sample account}{Drawings account}{Purchase account}
\mcq{101}{Discount received from Ali should be debited to:}{Discount received account}{Ali account}{Discount account}{Discount allowed}
\mcq{102}{Salaries paid to Zubair should be debited to:}{Salaries account}{Zubair account}{Cash account}{Bank account}
\mcq{103}{A table lost by fire should be debited to:}{Table account}{Furniture account}{Loss by fire account}{Purchase account}
\mcq{104}{Gave away charity in the form of goods should be debited to:}{Expense account}{Purchase account}{Cash account}{Charity account}
\mcq{105}{Credit signifies:}{Increase in liability}{Increase in revenue}{Increase in Capital}{All of above}
\mcq{106}{Purchase furniture from Saleem’s furniture at list price of Rs 500 subject to 5\% trade discount should be debited to:}{Purchase account}{Furniture account}{Furniture, trade discount}{Purchase, trade discount}
\mcq{107}{Purchased office typewriter from Muzammal should be credited to:}{Purchase account}{Office equipment account}{Muzammal account}{Cash account}
\mcq{108}{Sales of goods to X for cash should be debited to:}{Sales account}{Cash account}{X account}{None of these}
\mcq{109}{Payment of Rs 50 to Rahim as wages for repairing furniture should be debited to:}{Furniture account}{Repair account}{Wages account}{Rahim account}
\mcq{110}{Debit signifies:}{Increase in asset account}{Decrease in liability account}{Decrease in Capital account}{All of above}
\mcq{111}{A business transaction affects:}{At least one account}{At least two account}{Maximum two account}{Maximum three account}
\mcq{112}{Goods taken by the owner for personal use should be credited to:}{Sale account}{Drawing account}{Purchase account}{Cash account}
\mcq{113}{It is also called “King” of all books of accounts:}{Journal}{Ledger}{Trial balance}{Final account}
\mcq{114}{The process of recording transaction in ledger is called:}{Recording}{Journalizing}{Posting}{Transferring}
\mcq{115}{Accounts of assets normally have:}{Debit balance}{Credit balance}{Zero balance}{Both A,B}
\mcq{116}{Excess of credit over debit is called:}{Debit balance}{Credit balance}{Opening balance}{Closing balance}
\mcq{117}{A book in which transaction are recorded in chronological order is called:}{Ledger}{Trail balance}{Journal}{None of these}
\mcq{118}{A ledger is a book in which:}{Only Personal and Cash accounts are opened}{Only real accounts are opened}{Only nominal accounts are opened}{All real, nominal, personal accounts are opened}
\mcq{119}{A ledger is a book in which:}{Only Personal and Cash accounts are opened}{Only real accounts are opened}{Only nominal accounts are opened}{All real, nominal, personal accounts are opened}
\mcq{120}{If the two sides of an account are equal then account will show:}{Debit balance}{Zero balance}{Credit balance}{Debit \& Credit balance}
\mcq{121}{Accounts of Revenue normally have:}{Debit balance}{Credit balance}{Both debit and credit balance}{None of these}
\mcq{122}{Which one of the following normally has a debit balance:}{Capital account}{Creditor account}{Motor vehicles account}{Loan account}
\mcq{123}{Which of the following normally has a credit balance:}{Machinery account}{Purchase account}{Stock account}{Sales account}
\mcq{124}{Ledger is called the:}{Book of final entry}{Book of original entry}{Book of compound entry}{Book of single entry}
\mcq{125}{Left hand side of an account is called:}{Debit side}{Credit side}{Both of Debit \& Credit}{None of these}
\mcq{126}{The book in which transactions are finally recorded is called:}{Journal}{Ledger}{Trial balance}{Financial statement}
\mcq{127}{Which of the following is a liability:}{Stock}{Debtors}{Vehicles}{Creditors}
\mcq{128}{Double entry system is the system of:}{Recording transaction in books}{Recording in copies}{Presenting annual financial statement}{All of above}
\mcq{129}{The main purpose of trail balance is to:}{Check the equality of debits and credits}{Check the arithmetical accuracy of the double entry}{Provide information used in preparing final accounts}{All of above}
\mcq{130}{A trial balance is a:}{Balance sheet}{A rough draft of the final accounts}{List of the balanced in the ledger}{None of above}
\mcq{131}{Which of the following statement is true:}{Assets=Liabilities- Capital}{Capital=Assets + Liabilities}{Assets- Liabilities=Capital}{None of above}
\mcq{132}{The different stages of accounting cycle are:}{Documenting}{Recording}{Summarizing}{All of above}
\mcq{133}{Cash sales to Ahmad will be posted to debit side of:}{Cash account}{Sales account}{Ahmad account}{None of these}
\mcq{134}{The ruling of a trial balance is similar to that of:}{Ledger}{Final accounts}{Journal}{None of these}
\mcq{135}{Ledger records transactions in an:}{Chronological order}{Daily order}{Arithmetical order}{Analytical order}
\mcq{136}{Persons or institution are recorded in:}{Nominal accounts}{Real accounts}{Personal accounts}{Capital accounts}
\mcq{137}{Transactions relating to incomes and expenditures are recorded in:}{Real accounts}{Property accounts}{Nominal accounts}{Personal accounts}
\mcq{138}{Dividend received Rs;1500 from A will be posted to credit side of:}{Cash account}{Dividend account}{A,s account}{None of these}
\mcq{139}{Which of the following normally has a credit balance:}{Building account}{Purchases account}{Machinery account}{Sales account}
\mcq{140}{T account is a simplest form of a:}{Journal}{Ledger}{Trial balance}{Final account}
\mcq{141}{A bill of exchange is a method of payment on:}{Cash basis}{Credit basis}{Accommodation}{Reserve basis}
\mcq{142}{A bill of exchange is defined under Bills of exchange Act:}{1982}{1882}{1906}{1932}
\mcq{143}{A bill of exchange is an:}{Unconditional order}{A promise}{A request}{Conditional order}
\mcq{144}{A person who writes out the order to pay is called:}{Draw}{Acceptor}{Payee}{Drawer}
\mcq{145}{A person who is ordered to pay is called:}{Drawer}{Endorsee}{Creditor}{Draw}
\mcq{146}{When a drawer discounts the bill, he gets:}{Less than face value}{More than face value}{Equal to face value}{None of these}
\mcq{147}{It is a period of time after which a bill becomes payable:}{Grace days}{Tenor}{Usance}{Maturity}
\mcq{148}{The bills become payable at any time when they are presented by the holder are called:}{Trade bills}{Accommodation bills}{Term bills}{Demand bills}
\mcq{149}{In case of term bill extra three days are given to acceptor are called:}{Days of grace}{Days of tenor}{Days of Bill receivable}{Days of Bills payable}
\mcq{150}{A bill of exchange is drawn by:}{Debtor}{Creditor}{Endorser}{Endorsee}
\mcq{151}{If the drawer is in need of money and cannot wait till due date and receive the money from bank is called:}{Retirement of bill}{Endorsement of bill}{Discounting of bill}{Dishonor of bill}
\mcq{152}{When a drawer transfers a bill to his creditor for settlement of his own debts is called:}{Retirement of bill}{Endorsement of bill}{Discounting of bill}{Settlement of bill}
\mcq{153}{When a drawerhonors his acceptance he will be credited:}{Capital account}{Drawer account}{Bill payable account}{Cash account}
\mcq{154}{How many parties are involved in case of discounting of bill:}{Two}{Three}{Four}{None of these}
\mcq{155}{When a drawer discounts a bill, he debits:}{Bank account}{Interest account}{Drawee account}{Bill receivable account}
\mcq{156}{The person to whom the bill is transferred will become:}{Endorser}{Creditor}{Debtor}{Endorsee}
\mcq{157}{When a bill is endorsed, then the endorser has:}{Current liability}{Fixed liability}{Contingent liability}{None of those}
\mcq{158}{A bill from the point of view of a debtor is called:}{Bill acceptable}{Bill receivable}{Bill payable}{None of these}
\mcq{159}{In case of endorsement of bill, the endorser debits:}{Bill receivable account}{Bill payable account}{Cash account}{Creditor account}
\mcq{160}{“Bank for collection account” is a:}{Permanent account}{Fixed account}{Current account}{Temporary account}
\mcq{161}{When an acceptor refuses to pay the amount of bill to the holder of bill on its maturity is called:}{Honored bill}{Retired bill}{Dishonored bill}{Endorsed bill}
\mcq{162}{When a bill is dishonored, the holder of the bill goes to an official called:}{Inspector}{Notary public}{Drawer}{Manager}
\mcq{163}{When noting charges are paid at the time of dishonor of the bill, the drawee always debits:}{Interest account}{Discount account}{Noting charges account}{Drawer’s account}
\mcq{164}{In case of dishonor of bill, noting charges for dishonored bill are always paid by:}{Drawer of bill}{Notary public}{Acceptor of bill}{Holder of bill}
\mcq{165}{When the bill is dishonored acceptor’s account is debited in the books of the:}{Bank}{Endorsee}{Drawee}{Drawer}
\mcq{166}{Interest on renewal of a bill is an expense to the:}{Acceptor}{Drawer}{Endorsee}{Endorser}
\mcq{167}{Interest on renewal of a bill is an income to the:}{Drawee}{Debtor}{Acceptor}{Drawer}
\mcq{168}{When a drawee is declared insolvent, which account is to be debited in the books of drawer:}{Deficiency account}{Bad debts account}{Solvency account}{None of these}
\mcq{169}{In case of retirement of bill, the rebate or discount is a revenue for the:}{Holder}{Drawer}{Payee}{Acceptor}
\mcq{170}{Accommodation bills are drawn, accepted and endorsed:}{Without any consideration}{For quarter consideration}{For some consideration}{None of these}
\mcq{171}{Retiring a bill under rebate means:}{Payment of bill before due date}{Making payment for the bill after due date}{Discounting bill}{Dishonored bill}
\mcq{172}{Bill payable account is a:}{Nominal account}{Real account}{Personal account}{None of these}
\mcq{173}{A trader whose assets exceed from his liabilities is called:}{Debtor}{Creditor}{Solvent}{Insolvent}
\mcq{174}{Subsidiary books are called books of:}{Original entry}{Secondary entry}{Final entry}{Basic entry}
\mcq{175}{Which book is used to record all cash receipts and cash payments:}{Sale book}{Purchase book}{Cash book}{Petty cash book}
\mcq{176}{Which book is used to record all credit purchases:}{Purchases book}{Purchase returns book}{Cash book}{Debit note book}
\mcq{177}{All the transaction related to credit sales are recorded in:}{Sales returns book}{Sales book}{Cash book}{Credit note book}
\mcq{178}{The book which is used to record all goods returned by us to our supplier is called:}{Returns inwards book}{Purchases book}{Return outward book}{Sales book}
\mcq{179}{The book which is used to record all goods returned to us by our customers is called:}{Returns outwards book}{Returns inwards book}{Sales book}{Cash book}
\mcq{180}{All the acceptances received from the debtors are recorded in:}{Cash book}{Bills payable book}{Debtor’s book}{Bills receivable book}
\mcq{181}{All the bills given to the creditors are recorded in:}{Creditor’s book}{Bill payable book}{Bill receivable book}{Bank book}
\mcq{182}{Bill receivable book represents:}{Ledger}{Journal}{Trial balance}{All of above}
\mcq{183}{A book which is used to record small expenses is called:}{Expenses book}{Cash book}{Petty cash book}{Purchases book}
\mcq{184}{When the goods are returned to the suppliers, an intimation is sent to them is called:}{Currency note}{Invoice}{Debit note}{Credit note}
\mcq{185}{A document is sent to customer when he returns the goods is called:}{Promissory note}{Debit note}{Currency note}{Credit note}
\mcq{186}{The cash book is a book of:}{Compound entry}{Original entry}{Final entry}{Simple entry}
\mcq{187}{The periodical total of purchases book is posted to the:}{Debit of purchases account}{Credit of purchases account}{Credit of cash account}{Debit of sales account}
\mcq{188}{Both cash purchases as well as credit purchases are recorded in:}{Cash book}{Purchase book}{Purchases book}{None of these}
\mcq{189}{Sales returns book periodical total is posted to:}{Debit of returns inwards account}{Credit of returns inwards account}{Debit of returns outwards account}{Credit of returns outwards account}
\mcq{190}{The periodical total of purchases returns book is posted to:}{Debit of returns outwards account}{Credit of returns inwards account}{Debit of returns inwards account}{Credit of returns outwards account}
\mcq{191}{A document given by the seller to the buyer for credit sale of goods or assets is called:}{Cash memo}{Invoice}{Voucher}{None of these}
\mcq{192}{The book in which all cash transactions are primarily recorded is called:}{Pass book}{Cheque book}{Finance book}{Cash book}
\mcq{193}{When cash is received from debtor or customer, generally a receipt is issued to the debtor which is called:}{Payment voucher}{Receipt voucher}{Cash voucher}{None of these}
\mcq{194}{When money is paid to a creditor or supplier a receipt is obtained from him which is called:}{Cash voucher}{Cash payment}{Receipt voucher}{Payment voucher}
\mcq{195}{A document evidencing cash receipts and payments and forming the basis for making entries in the cash book is called:}{Receipt voucher}{Cash voucher}{Payment voucher}{Cash memo}
\mcq{196}{All cash purchases and cash sales are recorded in:}{Cash book}{Purchases book}{Sales book}{Purchases and sales book}
\mcq{197}{A cash account always shows:}{Debit balance}{Credit balance}{Both debit and credit balance}{Nil balance}
\mcq{198}{Credit purchases are recorded in:}{Purchases book}{Creditors book}{Cash book}{Purchases returns book}
\mcq{199}{The cash book is a:}{Subsidiary journal}{Ledger account}{Subsidiary journal and ledger}{Trial balance}
\mcq{200}{It is used to record only cash receipts and payments:}{Single column cash book}{Double column cash book}{Treble column cash book}{Sales book}
\mcq{201}{A double column cash book is used to record:}{Cash transactions only}{Cash and bank transactions}{Cash, bank and discount transactions}{None of these}
\mcq{202}{It is used to record cash, bank and discount transactions on debit and credit side:}{Single column cash book}{Petty cash book}{Double column cash book}{Treble column cash book}
\mcq{203}{An entry which is recorded on the both sides of cash book is called:}{Simple entry}{Compound entry}{Combined entry}{Contra entry}
\mcq{204}{The money can be deposited into bank by means of:}{Cash book}{Cheque book}{Pay-in-slip book}{Pass book}
\mcq{205}{When cheque received and deposited into bank, bank account must be:}{Credited}{Debited}{Both debited \& credited}{None of these}
\mcq{206}{When cash is withdrawn from the bank for the personal use of the owner, must be debited to:}{Cash account}{Bank account}{Drawings account}{Capital account}
\mcq{207}{Our cheque to Ali \& Bros, returned dishonored should be credited to:}{Cheque account}{Cash account}{Bank account}{Ali \& Bros account}
\mcq{208}{Ahmed’s cheque endorsed to Akbar should be debited to:}{Cash account}{Cheque account}{Akbar’s account}{Ahmad’s account}
\mcq{209}{Received a cheque from Bilal Rs 1900 in full settlement of his debt Rs2000 should be debited to:}{Cash account}{Discount account}{Bank account}{Botha,b}
\mcq{210}{Paid to Nadeem cash Rs. 5,000 and a cheque for Rs. 4800 in full settlement of his dues Rs10,000 should be credited to:}{Cash account}{Bank account}{Discount received account}{All of these}
\mcq{211}{Cash account is a:}{Personal account}{Nominal account}{Real account}{Both personal \& real account}
\mcq{212}{A statement which is used to reconcile the bank balance as per cash book with the balance as per pass book is called:}{Bank statement}{Income statement}{Bank reconciliation statement}{Financial statement}
\mcq{213}{Bank reconciliation statement is prepared by:}{Banker}{Customer’s accountant}{Auditors}{Manager}
\mcq{214}{The favorable balance of cash book bank column is:}{Credit balance}{Debit balance}{Both debit and credit balance}{None of these}
\mcq{215}{The favorable balance of bank statement is:}{Credit balance}{Debit balance}{Zero balance}{A,b,c}
\mcq{216}{Unfavorable balance mean:}{Credit balance in the pass book}{Credit balance in the cash book}{Debit balance in bank statement}{Both b,c}
\mcq{217}{Cheques issued to a creditor but not presented for payment are called:}{Uncredited cheques}{Uncollected cheque}{Unpresented cheques}{Dishonored cheques}
\mcq{218}{Cheques received from debtors but not collected by the bank are called:}{Dishonored cheques}{Uncredited cheques}{Unpresented cheques}{None of these}
\mcq{219}{The debit balances of cash and bank accounts in cash book are called:}{Income}{Liabilities}{Expenses}{Assets}
\mcq{220}{The credit balance of bank account in cash book is called:}{Expense}{Asset}{Income}{Liability}
\mcq{221}{Pass book is prepared by:}{Debtor}{Creditor}{Banker}{Customer}
\mcq{222}{A customer cheque returned dishonored is recorded in:}{Return outward book}{Return inward book}{Bill receivable book}{Cash book}
\mcq{223}{If receipt side of the cash book is under casted by Rs. 300, when overdraft as per bank statement is the starting point:}{Rs. 300 will be deducted}{Rs. 300 will be added}{Rs. 600 will be deducted}{Rs. 600 will be added}
\mcq{224}{Balance as per cash book overdraft Rs.15,000 interest on overdraft in pass book was Rs. 200 and interest on investment collected by bank Rs300 remained unrecorded in cash book, revised balance as per cash book will be:}{15,200}{15,300}{15,500}{15,900}
\mcq{225}{The balance of petty cash book is:}{A liability}{An asset}{A gain}{An expense}
\mcq{226}{A copy of the customer’s account in the banks ledger is called:}{Deposit statement}{Balance statement}{Bank statement}{Customer statement}
\mcq{227}{The term impress system is used in relation to:}{Cheque book}{Pay-in-slip book}{Cash book}{Petty cash book}
\mcq{228}{Any written evidence in support of a business transaction is called:}{Voucher}{Cash transaction}{Ledger}{Cash memo}
\mcq{229}{It is prepared to determine the gross profit or gross loss:}{Trading account}{Profit or loss account}{Balance sheet}{None of these}
\mcq{230}{It is prepared to determine the net profit or net loss:}{Trading account}{Profit or loss account}{Cash book}{Balance sheet}
\mcq{231}{Which of the following discloses the financial position of the business:}{Trading account}{Profit or loss account}{Profit or loss appropriation account}{Balance sheet}
\mcq{232}{Gross profit equals to:}{Net profit minus expenses}{Sales minus closing stock}{Purchases minus closing stock}{Sales minus cost of goods sold}
\mcq{233}{Cost of sales is equal to:}{Sales-purchase}{Purchases-return +closing stock}{Opening stock +purchases(net)- closing stock}{)Sales + opening stock-(purchases + closing stock)}
\mcq{234}{Net profit is equal to:}{Gross profit- expenses}{Sales –cost of goods sold}{Sales-expense}{Capital-expense}
\mcq{235}{The price of goods sold or service rendered to the customers is called:}{Sale}{Profit}{Expense}{Revenue}
\mcq{236}{Profit or loss appropriation account is not prepared in the case of:}{Partnership}{Joint stock company}{Sole trader-ship}{Pattern-ship at will}
\mcq{237}{Excise duty is a:}{Direct revenue}{Indirect revenue}{Direct expense}{Indirect expense}
\mcq{238}{The valuation of closing stock is at:}{Cost price}{Marker price}{Cost or market price whichever is lower}{Cost or market price whichever is higher}
\mcq{239}{Which account is a summary of direct expenses and direct revenue:}{Trading and profit or loss account}{Profit or loss account}{Balance sheet}{Trading account}
\mcq{240}{Net capital plus expenses is equal to:}{Purchases}{Cost of goods sold}{Capital}{Gross profit}
\mcq{241}{If the gross profit is Rs. 5,000 and the net profit is 35\% of the gross profit then the expenses must be:}{3250}{1250}{3750}{1750}
\mcq{242}{Position statement is similar to a:}{Trial balance}{Balance sheet}{Financial statement}{Bank reconciliation statement}
\mcq{243}{If sales are Rs. 12,000. Gross profit is 10\% of sales and net profit is 5\% of sales then the expenses will be:}{1200}{600}{1800}{2400}
\mcq{244}{Drawings are deducted from:}{Sales}{Income}{Capital}{Expenses}
\mcq{245}{Assets which have no physical existence are called:}{Tangible assets}{Fictitious assets}{Liquid assets}{Intangible assets}
\mcq{246}{An operating statement is similar to a:}{Balance sheet}{Bank Reconciliation statement}{Financial statement}{Profit or loss statement}
\mcq{247}{Stock in trade is a:}{Current asset}{Non-current asset}{Quick asset}{Intangible asset}
\mcq{248}{Assets which have no market value are called:}{Wasting assets}{Fictitious assets}{Intangible assets}{Tangible assets}
\mcq{249}{Net sales are equal to sales minus:}{Returns inwards}{Returns outwards}{Cost of goods sold}{Carriage on sales}
\mcq{250}{Goodwill, patent, copyright and trade mark are:}{Wasting assets}{Intangible assets}{Fictitious assets}{Liquid assets}
\mcq{251}{Debts which are repayable in the course of less than one year but more than one month are called:}{Quick liabilities}{Deferred liabilities}{Contingent liabilities}{Liquid liabilities}
\mcq{252}{Expenses related to sale of goods are shown in:}{Trading account}{Profit or loss account}{Balance sheet}{Sales account}
\mcq{253}{A balance sheet is a:}{Statement of income and expenditure}{Statement of debtors and creditors}{Financial statement of a business on a particular date}{Statement of profit earned by a business}
\mcq{254}{Closing stock is recorded in:}{Income statement only}{Trading account only}{Profit or loss account only}{Trading account and balance sheet}
\mcq{255}{A liability depends upon certain future event is called:}{Quick liability}{Deferred liability}{Contingent liability}{Current liability}
\mcq{256}{Marshaling of balance sheet means:}{Totaling of assets}{Totaling of liabilities}{Ordering of assets and liabilities}{Excess of assets over liabilities}
\mcq{257}{Bank overdraft is an example of:}{Current liability}{Liquid liability}{Quick liability}{All of these}
\mcq{258}{Income tax paid by a sole trader is shown on:}{Debit side of trading account}{Debit side of profit or loss account}{By way of deduction from capital in the balance sheet}{Credit side of trading account}
\mcq{259}{If the profit is 25\% of cost price then it is of what percentage of sale price:}{20\%}{25\%}{33\%}{66\%}
\mcq{260}{Gross profit is credited to:}{Trading account}{Profit or loss account}{Balance sheet}{None of these}
\mcq{261}{Net profit is transferred to:}{Creditor’s account}{Debtor’s account}{Drawing account}{Capital account}
\mcq{262}{It is a statement of assets, liabilities and owner’s equity on a particular date:}{Financial statement}{Balance sheet}{Bank reconciliation statement}{None of these}
\mcq{263}{Excess of assets over liabilities is called:}{Profit}{Income}{Expense}{Capital}
\mcq{264}{All the expenses connected with the management of the business are called:}{Office expenses}{Selling expenses}{Financial expenses}{Other expenses}
\mcq{265}{The method, under which the assets and liabilities are shown in Balance Sheet in the order of their permanence is called:}{Liquidity preference method}{Permanency preference method}{Mixed method}{None of these}
\mcq{266}{Permanence preference method is adopted by:}{Sole-trader-ship}{Partnership}{Joint stock company}{None of these}
\mcq{267}{The method under which assets and liabilities are shown in the order of their liquidity is called:}{Permanence preference method}{Liquidity preference method}{Mixed method}{Co-operative societies}
\mcq{268}{Liquidity preference method is usually adopted by:}{Joint stock company}{Sole trader-ship partnership}{Banks}{Insurance companies}
\mcq{269}{Assets, which have physical existence, are called:}{Tangible assets}{Intangible assets}{Quick assets}{Current assets}
\mcq{270}{The system of accounting in which accounting entries are made only when cash received or paid is known as:}{Cash system}{Accrual system}{Mercantile system}{Single system}
\mcq{271}{All those expenses which have become due but not paid are called:}{Accrued revenue}{Prepaid revenue}{Outstanding revenue}{Advance expense}
\mcq{272}{The accounting system in which accounting entries are made on the basis of amount having become due for payment or receipt, is called:}{Cash system}{Accrued or mercantile system}{Outstanding system}{None of these}
\mcq{273}{All those expenses which have not become due but paid in advance are called:}{Accrued expenses}{Outstanding expenses}{Prepaid expenses}{Payable expenses}
\mcq{274}{An income which has earned but not received is called:}{Unearned income}{Accrued revenue}{Revenue received in advance}{None of these}
\mcq{275}{The revenue that has not become due, but received in cash in current year is known as:}{Revenue received in advance}{Accrued revenue}{Unearned revenue}{Both A,C}
\mcq{276}{A prepaid expense is:}{An asset}{A liability}{An expense}{An income}
\mcq{277}{The accrued income or outstanding income will appear in the balance sheet as:}{An expense}{A liability}{An asset}{Both A,B}
\mcq{278}{Outstanding expense given in adjustment is called:}{An asset}{A gain}{An expense}{A liability}
\mcq{279}{A gradual decrease in the value of fixed assets is called:}{Reduction}{Revaluation}{Deprecation}{None of these}
\mcq{280}{Deprecation is:}{A liability}{A loss}{An expense}{Both B,C}
\mcq{281}{From business point of view, interest on capital is considered as:}{An income}{An expense}{A profit}{A liability}
\mcq{282}{Interest due but not received is an:}{Outstanding expense}{Accrued income}{Prepaid expense}{Unearned income}
\mcq{283}{Interest on drawings is:}{An asset}{A liability}{An expense}{A revenue}
\mcq{284}{Net profit is always:}{Equal to gross profit}{More than gross profit}{Less than gross profit}{None of these}
\mcq{285}{Wages paid for installing a machine should be debited to:}{Wage a/c}{Cash a/c}{Machinery a/c}{None of these}
\mcq{286}{Net loss is always:}{Equal to gross loss}{More than gross loss}{Less than gross loss}{None of these}
\mcq{287}{Income tax paid is a:}{Business expense}{Revenue for business}{Liability of business}{Personal expense}
\mcq{288}{Patent is an example of:}{Assets}{Liabilities}{Income}{Expense}
\mcq{289}{The amount becomes due from the customers is called:}{Allowance}{Loan}{Debts}{Doubtful debts}
\mcq{290}{The debts, which are irrecoverable from the debtors, are called:}{Debts}{Doubtful debts}{Bad debts}{Good debts}
\mcq{291}{Bad debts are business:}{Expenses}{Liabilities}{Assets}{Revenue}
\mcq{292}{Bad debts recovered from debtors should be credited to:}{Debtors a/c}{Bad debts a/c}{Bad debts recovered a/c}{Cash a/c}
\mcq{293}{The debts, the recovery or realization of which is uncertain are known as:}{Bad debts}{Doubtful debts}{Provision for doubtful debts}{Loan}
\mcq{294}{Provision for the doubtful debts is debited to:}{Debtors a/c}{Provision for bad debts a/c}{Profit or loss a/c}{Bad debts recovered a/c}
\mcq{295}{Business allow cash discount on making prompt payment to its:}{Creditors}{Debtors}{Partners}{Investors}
\mcq{296}{Business receive discount from its:}{Creditors}{Debtors}{Partners}{Owners}
\mcq{297}{Withdrawal of goods by the owner for his personal use must be credited to:}{Drawings a/c}{Goods a/c}{Purchases a/c}{Capital a/c}
\mcq{298}{If the closing stock appears in the balance it is taken only to the:}{Profit or loss a/c}{Trading a/c}{Balance sheet}{Final a/c}
\mcq{299}{The loss on the sale of old machinery is debited to:}{Machinery a/c}{Depreciation a/c}{Accumulated depreciation}{Profit and loss a/c}
\mcq{300}{Adjusting entries are made:}{During the accounting period}{At the end of accounting year}{In the beginning of the year}{At any time}
\mcq{301}{Provision for bad debts should appear in the balance sheet as:}{A fictitious asset}{Part of the share capital}{A deduction from the debtors}{An additional to the good will}
\mcq{302}{Bad debts are debited to:}{Debtor’s a/c}{Creditors a/c}{Provisions for bad debts a/c}{Bad debts a/c}
\mcq{303}{In the balance sheet pre-received income is shown on:}{Assets side}{Liabilities side}{Both assets \& liabilities sides}{None of these}
\mcq{304}{Provision for discount on creditors must be credited to:}{Creditors a/c}{Provision for discount on creditors a/c}{Profit and loss a/c}{Cash a/c}
\mcq{305}{Worksheet provides complete information for preparing:}{Income statement}{Position statement}{Financial statement}{Bank Reconciliation statement}
\mcq{306}{On way of avoiding errors in the permanent accounting records is to use a form which is called:}{Bank statement}{Trial balance}{Balance sheet}{Work sheet}
\mcq{307}{Adjusting entries affect:}{Balance sheet}{Income statement}{Work sheet}{All of those}
\mcq{308}{Which one of the following is not considered the permanent part of the accounting record:}{Journal}{Trial balance}{Ledger}{Work sheet}
\mcq{309}{The body of worksheet contains:}{Two pairs of money column}{Three pairs of money column}{Four pairs of money column}{Five pairs of money column}
\mcq{310}{In the preparation of the worksheet the amounts in the trial balance columns are taken directly from:}{Journal}{Ledger}{Trial balance}{Transaction}
\mcq{311}{Expired cost of fixed assets means:}{Appreciation}{Amortization}{Depreciation}{None of these}
\mcq{312}{A sheet which minimizes the chances of errors to the maximum possible extent is called:}{Balance sheet}{Work sheet}{Paper sheet}{Income statement}
\mcq{313}{Which one of the following is used to simplify work, to be done at the end of period:}{Balance sheet}{Income statement}{Bank statement}{Work sheet}
\mcq{314}{Financial statement is prepared is prepared from the amount given in:}{Trial Balance columns}{Adjustments columns}{Cash columns}{Adjusting trial balance columns}
\mcq{315}{The amount of assets, liabilities, capital and drawings are extended from adjusted trial balance to the:}{Balance sheet columns}{Income statement columns}{Cash columns}{Bank columns}
\mcq{316}{The excess of credit column over debit column in income statement is called:}{Net loss}{Net profit}{Cash balance}{Gross profit}
\mcq{317}{A statement, which is prepared to ascertain the net income or net loss for a specific accounting period is called:}{Position statement}{Statement of condition}{Financial statement}{Bank reconciliation statement}
\mcq{318}{A statement, which is prepared to known the financial position of a business on a particular date, is called:}{Balance sheet}{Statement of condition}{Position statement}{All of above}
\mcq{319}{The inflow of assets in return for services performed or products delivered during a period is:}{Asset}{Revenue}{Noncurrent asset}{None of these}
\mcq{320}{Cost incurred to generate revenue is called:}{Revenue}{Cost}{Expense}{Liability}
\mcq{321}{The amount by which the revenue for a particular period of time exceed from expenses incurred to generate them is called:}{Gross income}{Net income}{Revenue}{Expenses}
\mcq{322}{Sales are equal to:}{Cos of goods sold+profit}{Gross profit-cost of goods sold}{Cost of goods sold-gross profit}{None of these}
\mcq{323}{Cost of goods sold is equal to:}{Sales –purchases}{Purchase+closing stock-returns outwards}{Opening stock-closing stock+purchases+returns inwards}{Opening stock+purchases+return outwards-closing stock}
\mcq{324}{The gross amount of goods sold or services performed during an accounting period is known as:}{Purchase}{Sale}{Gross profit}{Net profit}
\mcq{325}{Net purchases are equal to:}{Sales-returns outwards}{Purchases+returnsPurchases inwards}{–returns outwards}{Purchases+returns outwards}
\mcq{326}{Net sales are equal to:}{Sale-returns outwards}{Sales +returns inwards}{Sales-returns inwards}{Purchases –returns inwards}
\mcq{327}{Which one of the following is a direct expense:}{Custom duty}{Carriage}{Wages}{All of above}
\mcq{328}{The margin that is available to cover the other expenses for a period and to yield net income is known as:}{Net profit}{Gross profit}{Net sales}{Cost of goods sold}
\mcq{329}{The expenses which are incurred for the generation of revenues from the sales of goods are called:}{Non-operating expenses}{Selling expense}{Operating expenses}{Administrative expenses}
\mcq{330}{All expenses connected with the office and its conduct are called:}{Selling expenses}{Marketing expenses}{Service expenses}{Administrative expenses}
\mcq{331}{Net sales-cost of goods sold-operating expenses=:}{Net purchases}{Cost of goods sold available for sale}{Net operating income}{Gross profit}
\mcq{332}{Income statement is prepared to ascertain the net result of the business concern under:}{British accounting system}{American accounting system}{Cash system of accounting}{None of these}
\mcq{333}{Which one of the following is a direct expense:}{Packing expenses}{Advertisement expenses}{Manufacturing expenses}{Travelling expenses}
\mcq{334}{Allowance for discount on account receivable is calculated on the amount of account receivable:}{Before deducting the allowance for doubtful debts}{Left after deducting the allowance for doubtful debts}{Before deducting the actual bad debts}{None of these}
\mcq{335}{Allowance for doubtful debts normally has a:}{Debit balance}{Credit balance}{Both debit \& credit balance}{None of these}
\mcq{336}{From the business point of view, interest on drawings is:}{An asset}{An expense}{A liability}{A revenue}
\mcq{337}{The amount invested by the owner in the business to produce revenue is known as:}{Income}{Asset}{Capital}{Liability}
\mcq{338}{It is the price of goods sold or services provided by a business to its customers:}{Asset}{Cost}{Capital}{Revenue}
\mcq{339}{The transactions, the effect of which is not exhausted with in the current accounting year are called:}{Revenue transactions}{Capital transactions}{Current transactions}{Monetary transactions}
\mcq{340}{Transactions, having short-term effects are known as:}{Revenue transactions}{Capital transactions}{Non-monetary transactions}{Paper- transactions}
\mcq{341}{An expenditure which is non-recurring and irregular, is called:}{Capital expenditure}{Revenue expenditure}{Short-term expenditure}{Current expenditure}
\mcq{342}{An expenditure, which is incurred to increase the profit earning capacity of a business concern, is called:}{Deferred expenditure}{Current expenditure}{Capital expenditure}{Recurring expenditure}
\mcq{343}{Wages paid for the construction of building is an example of:}{Revenue expenditure}{Capital expenditure}{Recurring expenditure}{Short-term expenditure}
\mcq{344}{An expenditure which is completely exhausted with in the current accounting period is known as:}{Deferred expenditure}{Revenue expenditure}{Future expenditure}{Non-recurring expenditure}
\mcq{345}{An expenditure, which is temporarily increase the profit-making capacity of the business is called:}{Deferred expenditure}{Capital expenditure}{Revenue expenditure}{Non-recurring expenditure}
\mcq{346}{Expenditure is a capital expenditure because:}{The amount involved is heavy}{It is the personal expenditure of the owner out of his capital}{It is intended to benefit the future period}{It is a recurring expenditure}
\mcq{347}{Expenditure is revenue expenditure because:}{It is intended to benefit the current period}{The amount involved is small}{It is deducted from the gross sale proceeds}{None of these}
\mcq{348}{Expenditure, which helps to maintain the business efficiency is called:}{Revenue expenditure}{Deferred expenditure}{Capital expenditure}{Future expenditure}
\mcq{349}{Which one of the following is appeared in the balance sheet:}{Revenue expenditure}{Capital expenditure}{Deferred expenditure}{Both b,c}
\mcq{350}{Depreciation of fixed assets used in the business is an example of:}{Capital expenditure}{Revenue expenditure}{Deferred expenditure}{None of these}
\mcq{351}{All revenue expenditures are taken to:}{Trading a/c}{Trading \& profit or loss a/c}{Profit or loss a/c}{Balance sheet}
\mcq{352}{An expenditure, incurred to improve the position of the business is known as:}{Deferred expenditure}{Revenue expenditure}{Capital expenditure}{Recurring expenditure}
\mcq{353}{Bad debts are:}{Deferred expenditure}{Revenue expenditure}{Capital expenditure}{None of these}
\mcq{354}{Octori duty paid on machinery, is an example of:}{Revenue expenditure}{Recurring expenditure}{Capital expenditure}{Both a,b}
\mcq{355}{Cost of redecorating a cinema hall is a:}{Capital expenditure}{Capital loss}{Revenue expenditure}{None of these}
\mcq{356}{A revenue expenditure, the benefit of which is not confined to one accounting year is called:}{Non- current expenditure}{Revenue expenditure}{Future expenditure}{Deferred expenditure}
\mcq{357}{An expenditure, which increases the utility or productive capacity of an asset is treated as:}{Revenue expenditure}{Capital expenditure}{Deferred expenditure}{None of these}
\mcq{358}{Heavy expenditure on advertisement for making a new product is a:}{Revenue expenditure}{Deferred expenditure}{Capital loss}{Non-recurring expenditure}
\mcq{359}{Distinction between capital and revenue items is important for the preparation:}{Balance sheet}{Trading and Profit or Loss a/c}{Bank Reconciliation statement}{Both a,b}
\mcq{360}{Distinction between capital and revenue items is important for the preparation:}{Balance sheet}{Trading and Profit or Loss a/c}{Bank Reconciliation statement}{Both a,b}
\mcq{361}{Capitalized expenditure are shown in:}{Trading a/c}{Profit or Loss a/c}{Income statement}{Balance sheet}
\mcq{362}{Preliminary expenses paid in the formation of a company is a:}{Capital expenditure}{Deferred expenditure}{Revenue expenditure}{Capital loss}
\mcq{363}{An expenditure incurred to keep the activities of a concern going on is:}{Capital expenditure}{Future expenditure}{Revenue expenditure}{None of these}
\mcq{364}{Receipts, which are non-recurring by nature, are called:}{Revenue receipts}{Current receipts}{Capital receipts}{Capital profit}
\mcq{365}{A receipt is a capital receipt because:}{The amount is small}{It relates to routine activity of the business}{It is received in the accounting year}{None of these}
\mcq{366}{A receipt is a capital receipt because:}{The amount is heavy}{It is credited to capital account}{It relates to fixed assets}{It is intended to benefit the future period}
\mcq{367}{Capital receipts are shown in the balance sheet on the:}{Asset side}{Liability side}{Debit side}{None of these}
\mcq{368}{Revenue receipts are shown in:}{Profit and Loss a/c(debit) side}{Profit and Loss a/c(credit) side}{Balance sheet(debit) side}{Trading a/c (debit) side}
\mcq{369}{A profit, which is earned on the sale of a fixed asset, is called:}{Revenue profit}{Capital profit}{Gross profit}{None of these}
\mcq{370}{The profit that is earned during the ordinary course of business is regarded as:}{Revenue profit}{Capital profit}{Deferred profit}{None of these}
\mcq{371}{Capital profit should be transferred to:}{Trading a/c}{Income statement}{Balance sheet}{Profit and Loss a/c}
\mcq{372}{The loss incurred on raising capital of a joint stock company is called as:}{Capital loss}{Revenue loss}{Capital reserve}{Normal loss}
\mcq{373}{Capital loss is shown in balance sheet on:}{Asset side}{Liability side}{Both a \& b}{Credit side}
\mcq{374}{Loss on sale of goods is an:}{Capital loss}{Revenue loss}{Revenue payment}{Deferred loss}
\mcq{375}{A loss is a revenue loss because:}{It is related to current assets}{It is incurred to decrease the tax liability}{It arises due to normal reasons}{None of these}
\mcq{376}{A loss is a capital loss because:}{It arises due to abnormal reasons}{It is another name given to drawings out of capital}{It relates to fixed asset}{All of these}
\mcq{377}{Rs.1000 paid as wages for erecting a machine should be debited to:}{Cash a/c}{Wages a/c}{Machinery a/c}{None of these}
\mcq{378}{A payment is a capital in nature when:}{It arises due to abnormal reasons}{The amount is heavy}{It relates to capital expenditure}{None of these}
\mcq{379}{Errors which affect one account can be:}{Errors of principle}{Errors of posting}{Error of omission}{None of these}
\mcq{380}{Errors of principle arises when:}{Any transaction is incorrectly recorded, either wholly or partially}{Any transaction is left wholly or partially}{Any transaction is affects one account}{Any transaction is recorded in fundamentally incorrect manner}
\mcq{381}{If a transaction has been completely omitted from the journal it will be considered:}{Error of commission}{Error of principle}{Error of omission}{None of these}
\mcq{382}{If goods purchased from Rahim for Rs.499, credited to Rehman’s account for Rs.499, this is an:}{Error of commission}{Error of omission}{Compensating error}{Error of principle}
\mcq{383}{Trade expenses of Rs.180 posted in the ledger as Rs. 810, it will be considered as:}{Error of principle}{Error of omission}{Error of casting}{Error of transposition}
\mcq{384}{When two or more than two errors occurred on the opposite side of the account and cancelled the effect of each other are called:}{Errors of omission}{Compensating errors}{Errors of commission}{Errors of principle}
\mcq{385}{The process of totaling the data at the end of period is called:}{Posting}{Casting}{Compensating}{Recording}
\mcq{386}{If a liability is recorded as income, it will be considered as:}{Errors of commission}{Error of omission}{Error of principle}{None of these}
\mcq{387}{Errors in casting of subsidiary books are called as:}{Error of omission}{Compensating errors}{Error of posting}{Clerical errors}
\mcq{388}{Errors of omission affects:}{One account}{Two accounts}{Three accounts}{None of these}
\mcq{389}{Transportation cost paid for the purchase of machinery must be debited to:}{Transportation cost account}{Purchases account}{Machinery account}{Cash account}
\mcq{390}{Errors of carry forward will affect:}{Personal accounts}{Nominal account}{Impersonal account}{None of these}
\mcq{391}{The credit purchases were wrongly recorded in sales book, the rectification of entry:}{Increase the net profit by double amount}{Decrease the net profit by double amount}{Decrease the net profit}{Increase the net profit}
\mcq{392}{If there is any error in Bank account it will affect:}{Trading and profit and loss account}{Profit and loss account}{Trading account}{Balance sheet}
\mcq{393}{If any expense omitted to be recorded it will:}{Overstate the profit}{Understate the profit}{Both a,b}{No effect on profit}
\mcq{394}{If any income omitted to be recorded it will:}{Overstate the profit}{Understate the profit}{Both overstate and understate}{Having no effect on profit}
\mcq{395}{A sale of Rs.1000 to Farid, was credited to his account, it will affect:}{Sales account}{Farid account}{Cash account}{Both a,b}
\mcq{396}{If the error committed in the capital account, it will affect:}{Trading account}{Profit and loss account}{Trading and profit and loss account}{Balance sheet}
\mcq{397}{Wages paid for the erection of a machine debited to wages account is an example of:}{Error of omission}{Error of commission}{Error of principle}{None of these}
\mcq{398}{Error which affects profit and loss account relates to:}{Nominal account}{Property account}{Personal account}{None of these}
\mcq{399}{Errors, which do not effects on profit calculation, will have an effect only on:}{Trial balance}{Balance sheet}{Profit or loss account}{Trading account}
\mcq{400}{Any difference in trial balance, is transferred to:}{Sales account}{Nominal account}{Purchase account}{Suspense account}
\mcq{401}{When balance of suspense account has debit balance it will be shown in balance sheet on:}{Liability side}{Capital side}{Asset side}{Credit side}
\mcq{402}{If amount for Rs.554 recovered from Ali, previously written off as bad debts we should:}{Debited to Bad debts recovered account}{Credited to Ali, account}{Debited to bad debts account}{Credited to Bad debts recovered account}
\mcq{403}{If the balance of suspense account is credit then it will be shown in balance sheet on:}{Assets side}{Asset and liability side}{Asset and capital side}{Capital and liability side}
\mcq{404}{Credit purchases from Ahmad costing Rs.180 were entered in the books for Rs.810 the rectifying entry is:}{Debit suspense account}{630Debit purchase account for Rs. 630 and credit Ahmad account for Rs. 630}{Debit Ahmad account for Rs. 630 and credit purchases account for 630}{None of these}
\mcq{405}{If sales return for Rs. 3,000 were incorrectly included in sales book, gross profit will be:}{Overstated by Rs.3,000}{Understated by Rs.6,000}{Overstated by Rs.6,000}{Understated by Rs.3,000}
\end{document}
